% Generated draft: requires mathematical and source review.
\documentclass[11pt]{article}
\usepackage[margin=25mm]{geometry}
\usepackage{fontspec}
\setmainfont{Latin Modern Roman}
\usepackage{amsmath,amssymb,mathtools,unicode-math}
\setmathfont{Latin Modern Math}
\usepackage{enumitem,needspace,xurl,hyperref}
\setlength{\parindent}{0pt}
\setlength{\parskip}{5pt}
\setlength{\emergencystretch}{3em}
\newcommand{\N}{\mathbb{N}}
\newcommand{\Z}{\mathbb{Z}}
\newcommand{\Q}{\mathbb{Q}}
\newcommand{\R}{\mathbb{R}}
\newcommand{\C}{\mathbb{C}}
\newcommand{\F}{\mathbb{F}}
\newcommand{\A}{\mathbb{A}}
\newcommand{\PP}{\mathbb{P}}
\newcommand{\E}{\mathbb{E}}
\newcommand{\eps}{\varepsilon}
\newcommand{\dd}{\,\mathrm{d}}
\begin{document}
\section*{UnsolvedMath ID: 30000218}
\textbf{Source problem code:} catalogue entry
\subsection{Definitions and mathematical statement}
Let $\mathcal H$ be the intended infinite-dimensional complex Hilbert space, $H_0$ the drift Hamiltonian, $\mu$ the control-coupling operator, and $\varepsilon(t)$ a scalar control. The verified source target concerns
\[
 \sqrt{-1}\,\partial_t\Psi(t)=\bigl(H_0-\varepsilon(t)\mu\bigr)\Psi(t),
 \qquad\Psi(0)=\Psi_0.
\]
For a chosen admissible control class $\mathcal U_T$ and a well-posed propagator $U_\varepsilon(T,0)$, the attainable states at time $T$ are
\[
 \mathcal R_T(\Psi_0)=\{U_\varepsilon(T,0)\Psi_0:
                                     \varepsilon\in\mathcal U_T\},
\]
and free-final-time reachability uses $\mathcal R(\Psi_0)=\bigcup_{T>0}\mathcal R_T(\Psi_0)$.

The question is to characterize these attainable final states in infinite dimension and give implementable criteria for the appropriate notions of controllability. For a specified target state space $\mathcal S$, exact controllability means equality with the relevant reachable set; approximate controllability means its closure equals $\mathcal S$ in a specified topology. Fixed-final-time and free-final-time versions differ. Under usual self-adjoint well-posedness hypotheses the norm is conserved, so a normalized state evolves on the unit sphere; if global phase is identified, the target is instead a projective state space. The source does not choose among these conventions.

No domains, spectral assumptions or boundedness hypotheses on $H_0,\mu$, allowable control regularity/amplitudes, final-time requirement, target topology or phase convention are supplied. In particular, the displayed evolution law alone does not establish a unitary propagator for every control or permit replacing infinite-dimensional reachability by a finite-dimensional Lie-algebra test. These inputs are needed before an exact characterization or computational criterion can be quantified.

The literature assessment checked on 21 August 2026 describes the original broad controllability prompt and says system-specific results do not settle it. The corrected statement restores the bilinear Schr\"odinger dynamics, but still does not fix the operators, controls or target set. Its open label therefore remains programme-level dated triage, not a theorem asserting controllability of every such equation.
\subsection{Short English statement}
Characterize the final states reachable by scalar controls in an infinite-dimensional bilinear Schr\"odinger system and give usable controllability tests. Exact versus approximate, time constraints, topology and global phase must be specified before the target becomes a single theorem.
\subsection{Sources}
\begin{itemize}
\item \emph{Optimal Control of Coupled Systems of PDE} (2006), Oberwolfach Reports, DOI 10.4171/owr/2005/18, pp. 995--1072, \url{https://doi.org/10.4171/owr/2005/18}; proposed year supplied as 2005.
\item The sole literature source in the assessment checked 21 August 2026 is the ``original Oberwolfach Report 18/2005,'' \url{https://doi.org/10.4171/owr/2005/18}. No particular system's theorem is supplied as a universal resolution.
\item Local UnsolvedMath registry, permanent ID $30000218$, exact source code \texttt{catalogue entry}, created 25 August 2026 and updated 16 September 2026. The clean statement is labelled corrected and verified against the official report, restoring the bilinear equation and initial condition.
\end{itemize}
\end{document}
